\section{引言}
\label{sec:introduction}

人工智能在历史上一直被工程化为一种很大程度上静态的制品。模型被训练，策略被选定，提示被撰写，架构被部署，由此产生的系统随后被评估，仿佛其智能是发布时固定的属性。这种观点便于基准测试和产品部署，但与大语言模型、工具使用型智能体、自动化程序合成和进化搜索等新兴系统日益不匹配。当代 AI 系统不仅仅是将输入映射到输出。它们能够检查执行轨迹、编写并执行代码、批判自身的推理、修改自身的提示、调用外部工具、存储情景记忆、生成新的训练任务，甚至对实现自身的软件提出修改建议。这些能力表明了一种从静态 AI 向\emph{自我进化 AI}的转变：即通过针对自身组件的反馈驱动循环来改进的系统。

本调查研究了这一转变。我们使用\emph{AI 自我进化}这一术语来指代一类方法，其中 AI 系统根据从与任务、用户、环境、验证器或其他智能体的交互中收集的反馈，修改自身的某些部分——包括其提示、记忆、工具、智能体代码、架构、数据课程或神经权重。该领域并非单一算法。它是一个汇聚区域，在这里，语言-智能体推理、AutoML 与神经架构搜索、进化计算、来自可验证反馈的强化学习、程序合成和开放式科学自动化开始共享概念。提示精炼循环（如 Self-Refine \cite{selfrefine2023}）、基于记忆的语言强化学习器（如 Reflexion \cite{reflexion2023}）、符号反向传播智能体优化器 \cite{agentsymboliclearning2024}、编写新智能体架构的元智能体 \cite{adas2025}、达尔文式自我修改编码智能体档案库 \cite{darwingodelmachine2025}、受哥德尔机器启发的运行时猴子补丁智能体 \cite{yin2024godel,schmidhuber2003}、进化编码系统（如 AlphaEvolve \cite{alphaevolve2025}）和零数据自博弈推理器 \cite{absolutezero2025}，在机制上各不相同。然而，它们可以在同一抽象下进行比较：系统观察性能，将反馈转化为修改提案，验证该提案，并有条件地保留修改后的系统。

这种汇聚的时机并非偶然。经典 AI 已经包含了这些要素：图灵对机器智能的行为框架 \cite{turing1950computing}、递归自我改进的思想实验、哥德尔机器理论 \cite{schmidhuber2003}、元学习、进化算法和 AutoML。在 2020 年代发生变化的是基础模型的可用性，这些模型能够操作构建现代 AI 系统所用的表示。提示是文本，工具描述是文本，代码是文本，基准报告是文本，堆栈追踪是文本，科学论文是文本，甚至许多架构决策也可以序列化为文本。因此，大语言模型充当了通用变异算子：它们可以提出新的提示、编辑 Python 类、设计测试、总结失败原因、合成训练样例，或解释先前尝试为何失败。当此类提案与自动评估器——单元测试、代码执行器、人类偏好模型、软件工程基准、数学验证器、网络任务成功标准或同行评审标准——相结合时，其结果就是一个实用的自我修改与选择的循环。

研究挑战在于理解这些循环何时构成真正的自我进化，而非普通的适应或离线优化。一个仅仅在部署后被工程师重新训练的系统，并非本文所用意义上的自我进化。一个在固定更新规则下被动接收更多数据的模型可能是在线学习，但不一定是自我进化。相反，一个保持权重不变但在每次失败后修改其记忆和提示的语言智能体，可以在其行为实际受控的表示层面进行自我进化。关键问题不在于被修改的组件是神经的、符号的还是过程式的。关键问题在于 AI 系统是否参与了针对其自身运行机制的修改的生成、评估和保留。

本章介绍本调查的背景、定义、边界和贡献。第~\ref{subsec:intro-background}~节追溯了从静态 AI 和图灵测试到哥德尔机器、AutoML、进化架构搜索、大语言模型和现代自我改进智能体的路径。第~\ref{subsec:intro-definition}~节提供了 AI 自我进化作为系统组件上反馈驱动修改过程的形式化定义。第~\ref{subsec:intro-related}~节将自我进化与持续学习、在线学习、自监督学习、元学习、AutoML、强化学习和传统智能体记忆进行区分。第~\ref{subsec:intro-contributions}~节总结了本调查的贡献：五层进化循环框架、对超过十九个开源和研究系统的分析、AutoML/NAS/LLM/进化计算/智能体研究的跨领域统一，以及该领域的研究者网络图谱。第~\ref{subsec:intro-structure}~节概述了本文的其余部分。

\subsection{背景：从静态 AI 到自我进化 AI}
\label{subsec:intro-background}

\subsubsection{早期人工智能中的静态系统假设}

对人工智能的最早公众想象已经与评估和变化的问题交织在一起。图灵 1950 年的提案将形而上学问题"机器能否思考？"替换为一种操作性模仿游戏，其中机器的智能通过交互来判断 \cite{turing1950computing}。该测试没有规定具体的学习机制，但它确立了一种持久的模式：智能将从任务协议下的行为中被推断出来。大多数后续基准测试继承了这一框架。系统被训练或编程，放置在接口之后，由其输出进行评判。协议可以是对话式的、数学的、感知的、策略性的或具身的，但被测试的对象通常在评估期间被视为固定的。

这种静态系统假设塑造了现代机器学习的组织方式。训练阶段产生参数 $\theta$；验证阶段选择超参数；测试阶段估计泛化能力。部署的模型随后被表示为函数 $f_{\theta}:\mathcal{X}\rightarrow\mathcal{Y}$，核心研究问题是如何从数据集 $D$ 中学习 $\theta$。即使是强调交互的强化学习，在报告基准分数时通常也会冻结已学习的策略。这种划分在科学上是有用的，因为它将搜索与评估分离开来，减少了泄漏，并使比较可复现。然而，它也将 AI 系统的概念缩小到最终制品，而不是产生并持续改变该制品的过程。

符号 AI 和专家系统在设计时之后也大多是静态的。工程师编码规则、分类法、产生式系统或搜索策略。知识库可以被更新，但更新通常由人类或狭窄的获取例程执行。专家系统不会重新设计其推理引擎、重写其知识获取工作流，或根据失败选择新的表示方案。系统的智能在于其手动指定的结构。机器学习后来用学习到的函数取代了手工编写的规则，但部署的制品仍然是一个快照。

静态假设因软件工程实践而得到强化。模型发布是版本化的、可审计的和受监控的。如果性能下降，开发者会收集数据、重新训练模型、修补代码或调整提示。系统通常没有被授予自行修改这些组件的权限。被改进的对象与改进它的人类过程之间的分离，正是 AI 自我进化开始侵蚀的主要边界。在自我进化系统中，改进过程的某些部分变成了系统自身运行的内在组成部分。

静态评估的历史重要性不应被低估。没有冻结的基准测试，就不可能衡量语音识别、计算机视觉、机器翻译、博弈或代码生成方面的进展。然而，同样的实践可能会遮蔽一类新型系统，这些系统的能力不仅存在于其当前策略中，还存在于其生成更好未来策略的能力中。进化智能体的静态快照类似于学习生物体的照片：有用但不完整。对于自我进化 AI，我们必须评估当前行为和随时间改变该行为的进化算子。

\subsubsection{自指与哥德尔机器的理想}

远在当前的语言智能体之前，研究人员就在思考智能机器是否能够改进自身的学习算法。Schmidhuber 的哥德尔机器是经典的形式化提案 \cite{schmidhuber2003}。它描述了一个自指系统，包含其自身软件的表示和一个证明搜索器。该机器仅在证明重写将根据形式化目标改善期望效用之后才重写自身。哥德尔机器的理论之美在于将自我修改视为一种理性行为。它不仅仅是变异自身；它要求证明变异是有益的。

原始哥德尔机器的理想远强于大多数现代系统。当前的智能体很少拥有其环境、效用函数或自身实现的完整形式化模型。它们通常无法证明某个代码补丁将改善长期效用。然而，这一思想引入了几项在当代自我进化中重新出现的原则。首先，系统必须能够访问对自身的某种描述。其次，自我修改必须相对于某个目标进行评估。第三，系统必须决定是否接受或拒绝候选修改。第四，改进过程是递归的：修改系统的机制本身也可能被修改。

现代系统将形式证明放宽为经验验证。它们不是证明某个修改全局地改善了效用，而是运行测试、基准、仿真或外部评估。例如，G{\"o}del Agent 明确引用了哥德尔机器的自指灵感，但通过 Python 猴子补丁和任务验证来实现运行时自我修改 \cite{yin2024godel}。Darwin G{\"o}del Machine 将自指雄心与达尔文式搜索相结合：基础模型提出对智能体的代码级修改，档案库保留在基准测试上有所改善的变体 \cite{darwingodelmachine2025}。这些系统并未提供 Schmidhuber 构造的数学保证。相反，它们以证明换取经验选择、可扩展性和与复杂软件环境的兼容性。

这种从证明到评估的转变是实用自我进化的决定性妥协之一。形式证明有吸引力但在开放领域中脆弱。经验评估有噪声但可扩展。本文调查的自我进化系统几乎总是使用后者。Reflexion 在失败的回合后存储语言反馈 \cite{reflexion2023}；Self-Refine 迭代地批判和修订输出 \cite{selfrefine2023}；Agent Symbolic Learning 在提示、工具和管道上产生语言梯度 \cite{agentsymboliclearning2024}；ADAS 评估新生成的智能体程序 \cite{adas2025}；AlphaEvolve 使用自动评估器选择程序变体 \cite{alphaevolve2025}；Absolute Zero 使用代码执行器作为自生成任务的验证器 \cite{absolutezero2025}。统一的模式不是携带证明的自我重写，而是携带反馈的自我修改。

哥德尔机器理论还澄清了为什么自我进化可能是不稳定的。如果系统可以修改自身的修改过程，它可能比仅限于固定更新规则的系统改进得更快。但同样的递归也可能放大错误、奖励黑客、不安全目标或分布漂移。改善短期基准的自我修改可能降低鲁棒性；增加某一任务成功率的提示修订可能移除安全约束；通过测试的代码补丁可能利用评估器。因此，自我进化既继承了自指的希望，也继承了其危险。本调查对循环、验证、回滚和治理的强调直接源于这种张力。

\subsubsection{从手工设计学习到 AutoML 和神经架构搜索}

第二条脉络来自学习系统设计的自动化。AutoML 提问将数据转化为模型的流水线本身是否可以被优化：数据预处理、特征工程、架构选择、超参数调优、训练调度和集成。神经架构搜索（NAS）将问题缩小到架构，通常使用强化学习、进化算法、可微松弛或权重共享。NASNet \cite{zoph2018learning}、Efficient NAS \cite{pham2018efficient}、DARTS \cite{liu2018darts} 和图像分类器的正则化进化 \cite{real2019regularized} 等方法表明，搜索可以发现与人类设计竞争的架构。

AutoML 和 NAS 对自我进化很重要，因为它们将智能从学习到的模型转移到了模型上的搜索过程。制品不再仅仅是一组权重；它还包括设计架构和训练配置的过程。然而，经典 NAS 通常是一个外部过程。控制器提出架构，训练循环评估它们，工程师部署优胜者。生成的模型通常在部署后不会继续搜索自身的架构。在这个意义上，AutoML 是自我进化的前身，但不总是其实例。

当搜索目标是智能体系统而非神经网络时，这种类比变得更加紧密。ADAS 将智能体系统的自动设计框架化为对实现智能体的程序的搜索 \cite{adas2025}。元智能体为新智能体架构编写代码，在任务上评估它们，并存储成功的设计。搜索空间是图灵完备的，因为智能体被表示为 Python 程序。这就是智能体软件的 NAS：系统不是选择卷积单元或 Transformer 块，而是搜索提示、工具调用、记忆策略、分解例程、反思过程和控制流。Agent Symbolic Learning 做了类似的转变，将智能体视为符号网络，其可学习元素是提示、工具和管道，而非可微权重 \cite{agentsymboliclearning2024}。该方法随后定义了类似于反向传播的语言损失、语言梯度和文本更新。

进化计算深化了这一脉络。NEAT 同时进化拓扑和权重 \cite{stanley2002evolving}；新颖性搜索论证了发散式探索可能比直接目标优化更高效 \cite{lehman2011abandoning}；MAP-Elites 和质量-多样性方法维护高性能多样化解决方案的档案库 \cite{mouret2015illuminating}；POET 在开放式过程中共同进化智能体和环境 \cite{wang2019poet}；AI-GAs 设想了生成架构、学习算法和环境的系统 \cite{clune2019ai}。现代自我进化智能体复用了这些思想，但用 LLM 引导的语义变异取代了随机或手工编码的变异算子。变异算子现在可以阅读代码、理解测试失败、编辑提示并提出结构化改进。

Darwin G{\"o}del Machine 是最清晰的合成。它维护智能体实现的档案库，采样父智能体，请求基础模型修改其代码，评估子代，并将改进的变体添加到档案库中 \cite{darwingodelmachine2025}。其报告的 SWE-bench 从 20.0\% 提升至 50.0\% 和 Polyglot 从 14.2\% 提升至 30.7\% 表明，智能体架构本身可以成为累积进化对象。AlphaEvolve 提供了工业级对应物：一个由 Gemini 驱动的进化编码智能体，将广泛探索与更深层次的推理和 MAP-Elites 式质量多样性相结合，报告了在矩阵乘法、数据中心调度、芯片设计和 LLM 训练内核方面的改进 \cite{alphaevolve2025}。这些系统明确化了 AutoML 所暗示的：设计过程可以成为被搜索、评估和改进的计算对象。

\subsubsection{大语言模型革命作为自我进化的使能因素}

大语言模型的兴起改变了自我进化的实际可行性。早期系统可以优化数值参数、离散架构或符号规则，但难以操作构成现代 AI 系统的异构制品：提示、工具 API、日志、源代码、测试结果、科学依据和自然语言反馈。GPT-4、Claude、Gemini 及相关前沿模型不仅是更强的文本预测器；它们是广泛胜任的软件和语言介导过程的操纵者。它们可以阅读堆栈追踪、推断 bug、生成补丁、产生单元测试、批判解决方案、重写提示、总结基准失败并提出新实验。

这种能力将语言变成了自我修改的通用接口。在 Self-Refine 中，模型生成输出，对其进行批判，并使用自身反馈进行精炼 \cite{selfrefine2023}。没有权重发生变化，但输出过程在单个回合内进化。Reflexion 增加了记忆：失败的尝试产生语言反思，作为后续尝试的条件 \cite{reflexion2023}。Agent Symbolic Learning 将类比推进了一步，在智能体计算图的节点上定义语言梯度 \cite{agentsymboliclearning2024}。RISE 使用强化学习将递归内省训练为一种习得的多轮行为 \cite{qu2024rise}。RAGEN 将轨迹分解为状态、思考、动作和奖励，同时识别出"回声陷阱"，即智能体通过重复自身而非推理来伪装修改 \cite{wang2025ragen}。ReVeal 将自我验证视为代码智能体中的一等习得能力，联合优化生成和验证 \cite{reveal2025}。

LLM 革命还使将代码作为进化基质的系统成为可能。SelfEvolve 是一个早期的代码进化框架，使用自生成的 API 知识和执行反馈迭代修复代码，改善了 DS-1000 和 HumanEval 性能 \cite{selfevolve2023}。ADAS 要求 LLM 元智能体编写新的智能体代码 \cite{adas2025}。G{\"o}del Agent 通过猴子补丁修改自身的运行时行为 \cite{yin2024godel}。Darwin G{\"o}del Machine 和 AlphaEvolve 使用基础模型作为智能体实现和算法代码上的语义变异算子 \cite{darwingodelmachine2025,alphaevolve2025}。AI Scientist 使用 LLM 智能体遍历更广泛的研究循环：想法生成、实验设计、代码执行、论文撰写和同行评审 \cite{aiscientist2024}。如果没有能够同时操纵自然语言和可执行制品的模型，这些系统是难以想象的。

重要的是，LLM 并未消除对评估的需求。它们使变异变得廉价且语义丰富，但它们也会产生幻觉、过拟合评分标准、利用漏洞，以及产生看似合理但无效的推理。因此，最强大的系统将 LLM 生成与外部反馈配对。代码执行器、测试套件、基准线束、环境奖励和人类或基于模型的评判器提供了选择压力。Absolute Zero 的零数据自博弈依赖于验证所提任务和解决方案的代码执行器 \cite{absolutezero2025}。AlphaEvolve 依赖于自动评估器 \cite{alphaevolve2025}。ReVeal 强调可靠的自我验证 \cite{reveal2025}。DGM 在软件工程基准上评估修改后的智能体 \cite{darwingodelmachine2025}。经验教训是 LLM 使自我进化成为可能，但验证器使其成为累积性的。

该领域的经验记录反映了这种对验证器的依赖。仅提示的自反馈可以改善许多任务，但也可能停滞或强化错误。基于记忆的反思在反馈具体且评估器可靠时可以提高通过率，但记忆可能无限增长。代码级自我修改可以发现新策略，但需要沙箱化、回滚和谨慎的接受标准。基于 RL 的自我改进可以将修订行为内化，但存在奖励黑客、推理坍塌或回声的风险。开放式档案库可以逃离局部最优，但代价高昂且难以治理。这些不是孤立的问题；它们是自我进化 AI 的核心设计权衡。

\subsubsection{从孤立技巧到自我改进系统}

最显眼的早期 LLM 自我改进方法是局部循环。单个答案被批判和修订。一段失败的代码样本被调试。一条反思被存储到记忆中。这些循环很重要，因为它们表明模型可以将自身输出用作改进材料。但它们不是该领域的终态。近期工作越来越多地将多个循环整合到跨回合、跨组件或跨种群改进的系统中。

一个有用的渐进过程从输出级精炼开始。Self-Refine 通过生成-批判-精炼迭代来改进候选输出 \cite{selfrefine2023}。Reflexion 通过存储先前失败的语言教训来改进未来尝试 \cite{reflexion2023}。SelfEvolve 通过结合自生成知识和解释器反馈来改进代码 \cite{selfevolve2023}。这些系统修改的是任务的上下文或制品，而不一定是智能体的源代码或权重。它们表明反馈可以被语言化表示并被复用。

下一步是组件级优化。Agent Symbolic Learning 将提示、工具和管道视为可学习的符号权重，允许部署的智能体在观察轨迹后调整内部组件 \cite{agentsymboliclearning2024}。DSPy 风格的提示编译、文本反向传播和多智能体提示进化属于同一家族。这里被改进的不仅是一个答案，而是智能体产生答案的过程。推理和学习之间的边界变得不那么清晰：智能体可以在执行任务的同时更新符号控制参数。

第三步是架构级和代码级进化。ADAS 使用元智能体编写新智能体 \cite{adas2025}。G{\"o}del Agent 修改运行时逻辑 \cite{yin2024godel}。DGM 通过档案库进化智能体实现 \cite{darwingodelmachine2025}。AlphaEvolve 大规模进化算法代码 \cite{alphaevolve2025}。这些系统将软件视为智能的基因组。基因组可由基础模型读取和编辑；评估决定哪些变体存活。这一视角使智能体设计成为累积性的且可能是开放式的。

第四步是权重级和课程级自我改进。RISE 通过多轮强化学习训练模型内省和修订其推理 \cite{qu2024rise}。RAGEN 将 RL 适配到多轮智能体，同时警告回声陷阱 \cite{wang2025ragen}。Absolute Zero 通过让模型在可验证奖励下提出任务并求解来消除外部训练数据 \cite{absolutezero2025}。ReVeal 联合优化代码生成和自我验证 \cite{reveal2025}。这些方法修改神经权重、任务课程或验证行为。它们表明自我进化不仅可以在提示或代码层面发生，还可以在模型的学习策略内部发生。

第五步是生命周期级自动化。AI Scientist 尝试自动化科学研究生命周期，从想法生成到实验执行和论文撰写 \cite{aiscientist2024}。生产级智能体框架如 AutoGen、LangGraph、Letta/MemGPT、OpenHands、SWE-Agent、CrewAI、MetaGPT、DSPy、OpenEvolve 及相关项目在其自我进化程度上差异很大，但它们共同展示了研究循环如何进入可用的工程生态系统。一些框架主要保持静态；另一些支持记忆进化、图修改、程序化智能体控制或进化代码搜索。因此，开源生态系统既是部署基底也是度量挑战：流行度、星标和热度不一定与真正的自我进化能力相关。

因此，背景促成了一个新的调查类别。AI 自我进化不等同于 LLM 提示，不归结为 AutoML，也不等同于强化学习。它是一种跨层设计模式，系统通过反馈修改自身。本章的其余部分将精确地定义这一模式。

\subsection{核心定义：形式化 AI 自我进化}
\label{subsec:intro-definition}

我们对\textbf{AI 自我进化}定义如下：

\begin{quote}
\textit{一个 AI 系统 $S$ 如果能够通过反馈驱动的循环修改自身的一个或多个运行组件——包括提示、记忆、工具、代码、架构、数据课程或权重——以改进在目标任务上的性能，并受决定候选修改是否被保留的接受机制约束，则该系统是自我进化的。}
\end{quote}

该定义强调五个要素：系统边界、可变组件、反馈、修改算子和保留。系统边界规定了什么算作"自我"。可变组件规定了什么可以被改变。反馈提供关于成功或失败的证据。修改算子提出变更。保留决定变更是否成为未来系统的一部分。没有反馈，变化是漂移而非进化。没有保留的修改，系统可能在单个回合内适应但不会累积改进。没有系统边界，每个外部人类更新都会算作自我进化，使该概念失去意义。

形式化地，设 AI 系统表示为
\begin{equation}
S_t = (\mathcal{C}_t, \mathcal{M}_t, \mathcal{E}, \mathcal{U}, \mathcal{G}),
\end{equation}
其中 $\mathcal{C}_t$ 表示时间 $t$ 的可变组件，$\mathcal{M}_t$ 表示记忆或状态，$\mathcal{E}$ 表示环境或任务分布，$\mathcal{U}$ 表示目标或效用准则，$\mathcal{G}$ 表示治理约束，如安全策略、沙箱限制或回滚规则。可变组件集合可以包括
\begin{equation}
\mathcal{C}_t = \{P_t, R_t, T_t, A_t, K_t, D_t, W_t\},
\end{equation}
其中 $P_t$ 是提示或指令，$R_t$ 是推理例程，$T_t$ 是工具和工具描述，$A_t$ 是智能体架构或控制流，$K_t$ 是源代码或可执行实现，$D_t$ 是数据或课程，$W_t$ 是神经权重。

在进化回合中，系统与任务 $x\sim\mathcal{E}$ 交互，产生轨迹 $\tau_t$，接收反馈 $F_t$，提出修改 $\Delta_t$，并在接受规则批准时更新其组件：
\begin{align}
\tau_t &= \mathrm{Run}(S_t, x), \\
F_t &= \mathrm{Feedback}(\tau_t, x, \mathcal{U}), \\
\Delta_t &= \mathrm{Modify}(S_t, F_t, \mathcal{M}_t), \\
S_{t+1} &= \mathrm{Accept}(S_t, \Delta_t, F_t, \mathcal{G}).
\end{align}
修改可以是提示编辑、记忆条目、工具更新、代码补丁、架构变更、生成的训练任务或权重更新。接受函数可以是确定性的（如仅在测试通过时保留补丁），或概率性的（如在进化种群中的选择）。它也可以包括人类批准，但系统必须在提出或评估变更方面做出实质性贡献。

性能改进目标可以写成
\begin{equation}
\mathbb{E}_{x\sim\mathcal{E}_{\mathrm{target}}}\left[ U(S_{t+1},x) \right]
\;>\;
\mathbb{E}_{x\sim\mathcal{E}_{\mathrm{target}}}\left[ U(S_t,x) \right] - \epsilon,
\end{equation}
其中 $\epsilon$ 允许噪声评估和探索。在严格的利用设置中，我们要求期望改进。然而，在开放式设置中，系统可能保留不立即支配父代但改善覆盖率、新颖性或未来可进化性的多样变体。对于基于档案库的系统，目标变为多维的：
\begin{equation}
\max_{S\in\mathcal{A}} \left( f(S),\; n(S),\; r(S),\; c(S)^{-1} \right),
\end{equation}
其中 $\mathcal{A}$ 是档案库，$f$ 是任务适应度，$n$ 是新颖性或多样性，$r$ 是鲁棒性，$c$ 是计算成本。这种形式捕获了质量-多样性方法，如 MAP-Elites \cite{mouret2015illuminating} 及其在 AlphaEvolve \cite{alphaevolve2025} 中的现代应用。

"其自身组件"这一表述应从操作性而非形而上学的角度来理解。自我进化的智能体不需要有意识、在法律意义上自主，或不受人类设计的约束。它只需具有一个循环，通过该循环，执行任务的系统由该系统或紧密耦合的元系统参与的修改过程来改变。例如，ADAS 使用元智能体设计目标级智能体 \cite{adas2025}。这是一个系统还是两个系统取决于边界。当元智能体、档案库、评估器和生成的智能体形成一个未来设计以自身搜索历史为条件的持久系统时，我们将其算作自我进化。类似地，Reflexion 的执行器、评估器和反思模型可以被视为一个复合系统，其记忆更新改变未来行为 \cite{reflexion2023}。

这种基于边界的观点避免了两个极端。一个极端是，每个训练过程都可以被称为自我进化，因为它改进了模型。这太宽泛了。标准监督学习通过外部指定的优化器在固定数据集上更新权重；模型不会提出对其自身学习过程或组件的修改。另一个极端是，只有重写自身源代码的系统才有资格。这太狭窄了。现代 AI 行为通常由提示、检索状态、记忆、工具和工作流图控制。如果这些组件被反馈驱动的循环修改，系统就在对行为有意义的层面上进化了。

因此，我们区分自我进化的级别：

\begin{enumerate}
    \item \textbf{输出级进化}：系统在任务回合内修订候选输出，如生成-批判-精炼循环。
    \item \textbf{记忆级进化}：系统存储改变未来行为的反馈、反思或经验。
    \item \textbf{提示/工具级进化}：系统编辑提示、工具描述或符号参数。
    \item \textbf{架构级进化}：系统改变控制流、模块组合或智能体拓扑。
    \item \textbf{代码级进化}：系统编写或修补实现自身或其后代的可执行源代码。
    \item \textbf{数据/课程级进化}：系统生成或选择塑造未来学习的任务。
    \item \textbf{权重级进化}：系统使用来自自生成或自中介经验的反馈更新神经参数。
    \item \textbf{种群级进化}：系统维护多个变体并随时间在其中进行选择。
\end{enumerate}

这些级别并非互斥的。ReVeal 结合了代码生成、自我验证、工具反馈和强化学习 \cite{reveal2025}。Absolute Zero 结合了任务生成、求解、可执行验证和权重更新 \cite{absolutezero2025}。DGM 结合了代码级修改与种群级档案管理 \cite{darwingodelmachine2025}。AI Scientist 结合了想法生成、实验代码、论文撰写和评审循环 \cite{aiscientist2024}。成熟的自我进化系统可能同时包含多个级别。

反馈是核心资源。它可以是内部的、外部的、符号的、标量的、文本的、可执行的或社会的。内部反馈包括自我批判和内省。外部反馈包括单元测试、环境奖励、基准分数、人类判断、同行评审分数和市场信号。符号反馈包括语言梯度和反思笔记。标量反馈包括奖励、通过/失败信号和性能指标。可执行反馈包括堆栈追踪、断言和验证器输出。社会反馈包括用户纠正和社区采纳。自我进化的可靠性在很大程度上取决于这种反馈的可靠性和一致性。弱评估器产生弱进化；可操纵的评估器选择的是操纵行为。

修改算子同样重要。在经典进化中，变异可能是随机的；在梯度下降中，更新方向由可微损失计算；在 AutoML 中，搜索算子是手工设计或控制器生成的。在现代自我进化 AI 中，LLM 通常实现修改算子：
\begin{equation}
\Delta_t = \mathrm{LLM}_{\phi}(\mathrm{serialize}(S_t), F_t, H_t),
\end{equation}
其中 $H_t$ 是历史、档案上下文或指令。该算子很强大，因为它可以对异构制品进行语义一致的编辑。但它也很危险，因为它是随机的、对提示敏感的，并且能够引入隐性故障。修补 bug 的同一个 LLM 可能移除安全防护；改进基准的同一个智能体可能过拟合到其公开测试。

保留机制决定了改进是否累积。一个天真的系统可能在每次自我批判后覆盖自身，产生漂移。更鲁棒的系统在验证任务上评估候选变更，与基线比较，存储回滚点，并保持多样性。Reflexion 的记忆缓冲区保留反思但必须管理上下文限制 \cite{reflexion2023}。G{\"o}del Agent 仅在验证后接受修改，并可以回滚失败的补丁 \cite{yin2024godel}。DGM 将变体存储在档案库中而非单一的可变谱系中 \cite{darwingodelmachine2025}。AlphaEvolve 使用质量-多样性选择在行为生态位中保留精英 \cite{alphaevolve2025}。保留是学习变为进化的关键所在。

最后，自我进化必须随时间评估。一次更新后的静态基准分数是不够的。我们必须追问系统是单调改进还是周期性波动，增益是否跨任务迁移，是否维持安全约束，是否会坍塌为重复行为，是否过拟合到评估器，以及进化成本是否由收益证明合理。对于进化轨迹 $S_0,S_1,\ldots,S_T$，相关指标包括累积改进，
\begin{equation}
\Delta_T = U(S_T)-U(S_0),
\end{equation}
每步效率，
\begin{equation}
\eta_T = \frac{U(S_T)-U(S_0)}{\sum_{t=0}^{T-1} c_t},
\end{equation}
以及保留分布下的鲁棒性，
\begin{equation}
\rho_T = \mathbb{E}_{x\sim\mathcal{E}_{\mathrm{heldout}}}[U(S_T,x)] - \mathbb{E}_{x\sim\mathcal{E}_{\mathrm{train}}}[U(S_T,x)].
\end{equation}
这些指标预示了本调查的评估章节，在那里我们论证自我进化系统需要衡量轨迹而非仅仅终点的基准。

\subsection{相关概念与边界}
\label{subsec:intro-related}

AI 自我进化与若干成熟领域存在交集，但并不等同于其中任何一个。本节划定边界，因为"自我改进"一词常被宽泛使用。一个系统可能在部署后有所改进、实现在线学习、适应用户偏好、调整超参数或在自监督损失上训练，却不满足自我进化的更强标准。反过来，一个系统的神经网络权重可能保持不变，但只要反馈驱动的提示词、记忆、工具、代码或架构变更改变了其未来行为，该系统仍可能是自我进化的。

\paragraph{持续学习。} 持续学习研究的是在不发生灾难性遗忘的情况下，从一系列任务或数据分布中进行学习的模型。其核心问题包括稳定性-可塑性权衡、经验回放、正则化、任务推断和终身适应。持续学习者随时间更新参数，但更新规则通常由设计者固定。模型不一定检查失败、提出新工具、重写自身架构或选择自己的课程。当系统参与修改自身的学习机制或任务接口，而不仅仅是在固定优化器中流过更多数据时，持续学习才成为自我进化。

\paragraph{在线学习。} 在线学习在样本到达时更新预测。通常以与某个比较器的遗憾界来形式化。学习者接收输入、预测、观察损失并更新。这是反馈驱动的，但反馈通常由固定算法使用，如随机梯度下降、乘法权重或赌博机更新。自我进化需要一个额外层次：系统可以改变表示、提示词、代码、架构甚至更新规则本身。RISE 和 RAGEN 比一般在线学习更接近自我进化，因为它们学习多轮自省或智能体策略，但即使如此，仍需分析所学行为是否改变了未来的自我修改能力 \cite{qu2024rise,wang2025ragen}。

\paragraph{自监督学习。} 自监督学习从数据本身构建标签，如掩码语言建模、对比学习或下一词元预测。它在"自"的意义上是自动推导监督信号，而非人工标注。它不一定是自我进化的，因为模型不会通过关于任务性能的反馈来修改自身的操作组件。Absolute Zero 是一个有启发性的对比：它生成任务和解决方案，用代码执行器验证，并通过强化学习更新 \cite{absolutezero2025}。这超越了自监督，因为模型在可验证奖励下塑造了自己的课程。

\paragraph{元学习。} 元学习训练跨任务"学习如何学习"的系统。它与自我进化密切相关，因为它关注适应机制。然而，许多元学习仍然是离线的：元学习器在任务分布上训练，然后使用固定的内循环过程进行适应。自我进化强调部署后或系统内部对组件的修改，通常包括可微元学习无法捕获的符号和代码产物。ADAS 可以解释为对智能体设计的元学习，但它也是自我进化的，因为搜索历史和生成的智能体制约着未来的设计 \cite{adas2025}。

\paragraph{AutoML 与 NAS。} AutoML 自动化模型设计、训练和选择。NAS 搜索架构空间。这些领域是直接先驱，但大多数 AutoML 流水线是外部优化器。人类调用优化器，等待候选模型，然后部署。AI 自我进化将该优化器的一部分内化到系统生命周期中。当元智能体设计新智能体、当智能体实现的存档不断累积、或当编码系统在评估器下变异自身算法时，AutoML 就变成了智能体化的和自我指涉的。

\paragraph{强化学习。} 强化学习通过奖励信号优化策略。它是自我进化的基础，特别是在 RISE、RAGEN、Absolute Zero 和 ReVeal 中 \cite{qu2024rise,wang2025ragen,absolutezero2025,reveal2025}。然而，仅 RL 并不意味着自我进化。由 PPO 在环境奖励上训练的策略可能变得更胜任，但不一定修改其提示词、工具、代码、架构或课程。当策略控制或参与修改过程时，如提出任务、生成测试、修正推理或更改智能体软件，自我进化才出现。

\paragraph{RLHF 与偏好优化。} 从人类反馈的强化学习将模型对齐到偏好 \cite{christiano2017deep,ziegler2019fine}。它使用反馈，但过程通常由外部编排：人类标注数据，工程师训练奖励模型，训练流水线更新权重。Constitutional AI 和自我批评方法通过在宪法下使用模型生成的批评和修订，更接近自我进化 \cite{constitutionalai2022}。尽管如此，区别在于部署后的系统是否能在反馈下继续修改自身组件，而不是仅仅用合成偏好训练一次。

\paragraph{智能体记忆与个性化。} 许多智能体框架存储记忆、用户偏好或对话摘要。如果记忆是从反馈中生成并保留以改进未来行为，那么记忆可以是自我进化的。Letta/MemGPT 式的记忆系统因此是部分自我进化机制。但仅有记忆是不够的。记录转录但不用于未来修改的日志系统不是自我进化的。存储用户偏好的推荐系统可能是自适应的，但如果其更新规则和任务接口是固定的，它属于更窄的类别。

\paragraph{程序合成与自动修复。} 程序合成从规约生成代码；自动程序修复根据测试修补缺陷。现代 LLM 代码智能体模糊了这个边界。SelfEvolve、ReVeal、AlphaEvolve 和 DGM 都将程序合成或修复作为自我进化的一部分 \cite{selfevolve2023,reveal2025,alphaevolve2025,darwingodelmachine2025}。区别特征是递归使用：代码生成不仅用于解决外部编程任务，还用于改进系统自身未来的问题求解过程。

\begin{table}[t]
\centering
{\TableTight
\begin{tabularx}{\textwidth}{@{}L{0.16\textwidth}L{0.20\textwidth}L{0.23\textwidth}Y@{}}
\toprule
\textbf{概念} & \textbf{主要变更对象} & \textbf{典型反馈} & \textbf{与 AI 自我进化的区别} \\
\midrule
持续学习 & 任务流上的神经参数 & 新数据、任务标签、回放损失 & 通常使用固定更新规则；可能不修改提示词、工具、代码、架构或课程。 \\
在线学习 & 每个样本后的预测器状态 & 即时损失或奖励 & 关注固定算法下的遗憾；自我进化可以改变算法机制本身。 \\
自监督学习 & 表示和权重 & 自动生成的标签 & 监督是自动的，但系统不一定评估和修改自身的操作组件。 \\
元学习 & 适应过程或初始化 & 训练任务分布 & 通常离线学习；自我进化强调运行时反馈循环和持久的系统修改。 \\
AutoML/NAS & 模型流水线或架构 & 验证性能 & 通常是外部设计优化器；当集成到智能体/系统生命周期中时变为自我进化。 \\
强化学习 & 策略参数 & 环境奖励 & RL 是优化基底；自我进化额外关注组件和过程的自我修改。 \\
RLHF/偏好调优 & 模型行为和对齐 & 人类或 AI 偏好标签 & 通常外部编排；自我进化要求持续的系统介导修改。 \\
智能体记忆 & 情景或语义上下文 & 用户交互、反思 & 若记忆源自反馈且行为活跃则为部分自我进化；否则仅为日志/个性化。 \\
程序合成/修复 & 任务特定代码 & 规约、测试、错误 & 当生成或修复的代码改变未来智能体/系统时为自我进化，而非仅一次性输出。 \\
\bottomrule
\end{tabularx}}
\caption{AI 自我进化与相邻概念的关系。边界是可渗透的：许多现代系统组合了多行内容，但自我进化特别要求反馈驱动的对系统自身操作组件的修改。}
\label{tab:related-concepts}
\end{table}

表~\ref{tab:related-concepts} 中的比较提示了一条一般规则。如果系统仅仅改变其输出，那是适应。如果它改变了影响未来输出的内部状态，那可能是学习。如果它在反馈和保留机制下，改变了学习、推理、行动或设计后继者的组件或过程，那就是自我进化。这条规则同时涵盖了轻量级的提示词/记忆方法和重量级的代码/权重/种群方法。

若干边界案例值得关注。首先，使用思维链提示的 LLM 如果每个响应是独立的，则不是自我进化。如果思维链被批评、修订、存储并复用，它可能成为输出级或记忆级的自我进化。其次，人类在观察失败后编辑智能体提示词不是 AI 自我进化，尽管它可能是人在回路进化系统的一部分。如果智能体提出编辑建议而人类仅审批，则系统是部分自我进化的。第三，针对基准调优的模型发布后不再是自我进化的，但训练组织可能实现外部进化流水线。本调查聚焦于其循环在技术上被表示且可作为 AI 系统一部分被分析的系统。

另一个边界涉及自主性。自我进化系统不必完全自主。人类审批、安全过滤器、沙箱和治理通常是可取的。事实上，不受限制的自主性是安全风险。因此，我们通过参与修改来定义自我进化，而非通过缺乏人类监督。一个为自身起草代码补丁、在沙箱中评估并在部署前请求审批的系统，仍然展现自我进化结构。接受函数将人类作为治理 $\mathcal{G}$ 的一部分纳入。

最后一个边界涉及改进。进化系统可能倒退。一次失败的修改不会使系统丧失资格，正如生物谱系可以积累有害突变。重要的是系统包含一个反馈驱动的选择过程，旨在随时间改进性能或多样性。评估应测量轨迹，包括失败，而非假设单调进展。

\subsection{本调查的贡献}
\label{subsec:intro-contributions}

本调查做出四项主要贡献。第一，我们提出了\textbf{五层进化循环}框架，按反馈转化为保留变更的机制来组织该领域。第二，我们分析了超过十九个研究和开源系统，将真正的自我进化机制与静态智能体编排和炒作驱动的采用区分开来。第三，我们在统一的正式抽象下整合了 AutoML/NAS、LLM 自我改进、进化计算、强化学习、程序合成和智能体框架等领域的工作。第四，我们绘制了产生该领域的研究者和机构网络，突出了进化计算谱系、语言智能体推理谱系和 LLM 自我改进谱系的汇合。

\subsubsection{贡献一：五层进化循环框架}

第一项贡献是一个按执行修改和选择的循环来分类自我进化 AI 的框架。我们使用五个循环，因为它们捕获了文献中的主导机制，同时保持足够的通用性以比较混合系统。

\paragraph{循环 I：反射性输出与记忆进化。} 系统通过批评输出、存储反馈和调节未来尝试来改进。Self-Refine 和 Reflexion 是代表性工作 \cite{selfrefine2023,reflexion2023}。被修改的组件通常是当前输出、情景记忆或反思缓冲区。评估器可以是同一模型、外部奖励模型、测试或环境。该循环成本低、易于部署，但受限于上下文长度、自我批评质量和评估器可靠性。

\paragraph{循环 II：符号组件进化。} 系统将提示词、工具、工作流和流水线视为可学习的符号参数。Agent Symbolic Learning 是核心示例：它为提示词、工具和流水线节点定义语言损失、语言梯度和优化器 \cite{agentsymboliclearning2024}。DSPy 式提示词编译和文本反向传播也属于此类。该循环桥接了基于梯度的学习和智能体工程。它使智能体的设计面可优化，而无需神经权重更新。

\paragraph{循环 III：验证驱动的代码进化。} 系统生成、测试、修复和保留代码。SelfEvolve 使用自生成知识加解释器错误 \cite{selfevolve2023}；ReVeal 协同优化生成和自我验证 \cite{reveal2025}；AlphaEvolve 使用自动化评估器选择算法程序 \cite{alphaevolve2025}。该循环之所以强大，是因为代码可以被机械地检查。其局限在于领域依赖性：可靠的验证器在编程、数学和仿真中比在开放式社会任务中更容易实现。

\paragraph{循环 IV：架构与智能体设计进化。} 系统在智能体架构、控制流、工具使用和自我修改代码上进行搜索。ADAS、G{\"o}del Agent 和 DGM 是代表性工作 \cite{adas2025,yin2024godel,darwingodelmachine2025}。被修改的组件是智能体本身。该循环最接近自我重设计的直觉概念。它可以发现人类未曾预见的策略，但也带来成本、安全性和可复现性挑战。

\paragraph{循环 V：课程、权重与种群进化。} 系统改变数据分布、训练课程、神经权重或种群存档。RISE、RAGEN、Absolute Zero、DGM 和 AlphaEvolve 都实例化了该循环的部分内容 \cite{qu2024rise,wang2025ragen,absolutezero2025,darwingodelmachine2025,alphaevolve2025}。该循环可以产生比纯提示词方法更深层的行为变化，但需要仔细的奖励设计和崩溃防护。Absolute Zero 作为任务提出者和求解者的双重角色特别重要，因为它消除了人工策划训练数据的假设。

这些循环可以概括为
\begin{equation}
\mathrm{SelfEvolution} = \mathrm{Observe} \rightarrow \mathrm{Interpret} \rightarrow \mathrm{Modify} \rightarrow \mathrm{Verify} \rightarrow \mathrm{Retain},
\end{equation}
每个循环在修改的表示和使用的验证器上有所不同。该框架不是论文的分类法，而是机制的分类法。单篇论文可以实例化多个循环。例如，DGM 结合了架构/代码进化与种群存档；ReVeal 结合了验证驱动的代码进化与权重级 RL；AI Scientist 结合了想法、代码、稿件和评审循环。

\begin{table}[t]
\centering
{\TableTight
\begin{tabularx}{\textwidth}{@{}L{0.18\textwidth}L{0.19\textwidth}L{0.25\textwidth}Y@{}}
\toprule
\textbf{循环} & \textbf{可变组件} & \textbf{代表系统} & \textbf{典型失败模式} \\
 \midrule
反射性输出与记忆 & 输出、反思、记忆缓冲区 & Self-Refine, Reflexion, ReAct 式智能体 & 泛化批评、记忆膨胀、上下文过拟合、局部最优。 \\
符号组件 & 提示词、工具、流水线、工作流节点 & Agent Symbolic Learning, DSPy, EvoMAC 式文本优化 & 提示词过拟合、不稳定语言梯度、高 LLM 调用成本。 \\
验证驱动代码 & 候选代码、测试、修复补丁 & SelfEvolve, ReVeal, AlphaEvolve, 代码智能体 & 测试过拟合、沙箱逃逸风险、验证器不完整性。 \\
架构与智能体设计 & 智能体代码、控制流、工具拓扑 & ADAS, G{\"o}del Agent, DGM & 不安全的自我修改、不可复现的搜索、昂贵的评估。 \\
课程、权重与种群 & 训练任务、策略权重、存档 & RISE, RAGEN, Absolute Zero, DGM, AlphaEvolve & 奖励博弈、回声陷阱、模式崩溃、计算成本。 \\
\bottomrule
\end{tabularx}}
\caption{本调查贯穿使用的五层进化循环框架。循环分类的是机制而非互斥系统。}
\label{tab:five-loops}
\end{table}

\subsubsection{贡献二：研究系统与开源项目分析}

第二项贡献是对超过十九个系统的面向证据的分析。研究核心包括 Self-Refine、Reflexion、Agent Symbolic Learning、ADAS、G{\"o}del Agent、Darwin G{\"o}del Machine、SelfEvolve、RISE、RAGEN、Absolute Zero、ReVeal、AlphaEvolve、FunSearch、AI Scientist、STaR、Constitutional AI、Tree of Thoughts、ReAct、SWE-agent \cite{sweagent2024} 及相关软件工程智能体。开源和框架生态包括 DSPy、AutoGen、LangGraph、Letta/MemGPT、CrewAI、MetaGPT、OpenHands、OpenEvolve、AutoGPT、smolagents 以及上述论文的项目特定实现。

我们的分析强调机制而非流行度。本调查使用的星标分析报告显示，GitHub 星标数可能是技术质量的糟糕代理。AutoGPT 积累了大量关注，但其星标-质量评分估计远低于较小的研究导向项目。OpenEvolve、DSPy、SWE-Agent、OpenHands、LangGraph、AutoGen、CrewAI 和 MetaGPT 展现了不同的传播模式：学术型、稳步型、炒作驱动型或社区驱动型。这很重要，因为自我进化是一个技术属性，而非营销标签。一个项目可能因为易于演示而流行，而一个更安静的项目可能包含更重要的自我进化机制。

因此，我们从可变组件、反馈来源、保留机制、评估基准、开放性、安全控制和迁移等维度检查每个系统。Reflexion 为记忆级语言反馈提供了有力证据，因为它在外部评估下报告了显著的 HumanEval 和交互式任务增益 \cite{reflexion2023}。Self-Refine 证明同模型批评可以在多样任务上平均改进约 20\%，但也说明了无外部验证的自反馈局限 \cite{selfrefine2023}。Agent Symbolic Learning 通过优化符号智能体组件，报告了 HotPotQA、MATH、HumanEval、软件开发可执行性和创意写作方面的改进 \cite{agentsymboliclearning2024}。DGM 通过开放式代码级进化报告了显著的软件工程改进 \cite{darwingodelmachine2025}。AlphaEvolve 在自动化评估下报告了算法和基础设施改进 \cite{alphaevolve2025}。Absolute Zero 证明自生成课程可以支持零数据推理改进 \cite{absolutezero2025}。

开源生态对理解采用情况同样重要。AutoGen 支持多智能体对话和人在回路工作流；LangGraph 提供可以编程方式修改或编排的状态图；Letta/MemGPT 强调长期记忆；DSPy 编译提示词和演示；SWE-Agent 和 OpenHands 将软件工程智能体操作化；OpenEvolve 为社区改编进化编程思想。CrewAI 和 MetaGPT 普及了基于角色的智能体编排，但许多部署仍然是静态的，除非与反馈驱动的修改相结合。区分静态编排与自我进化是本调查的实践目标之一。

\subsubsection{贡献三：跨领域统一}

第三项贡献是跨领域统一。同一抽象循环在不同领域以不同名称出现。在 AutoML 中是架构搜索。在进化计算中是变异和选择。在 LLM 提示词中是批评和精炼。在强化学习中是策略改进。在程序合成中是生成与测试。在软件工程中是自动修复。在智能体框架中是记忆和工作流适应。在 AI for Science 中是假设生成、实验和评审。

我们使用共享元组来形式化这些领域：表示、变异算子、反馈通道、保留规则和治理约束。对于 NAS，表示是神经架构，变异是架构变异或控制器采样，反馈是验证准确率，保留是选择，治理是计算预算。对于 Reflexion，表示是记忆，变异是语言反思，反馈是评估器奖励，保留是记忆更新，治理是上下文管理。对于 DGM，表示是智能体源代码，变异是 LLM 补丁，反馈是基准性能，保留是存档插入，治理是沙箱和选择。对于 Absolute Zero，表示包括任务课程和权重，变异是自生成任务和解决方案，反馈是代码执行，保留是 RL 更新，治理是验证器设计。

这种统一揭示了共享瓶颈。搜索成本出现在 NAS、ADAS、DGM 和 AlphaEvolve 中。评估器可靠性出现在 Self-Refine、Reflexion、ReVeal、AI Scientist 和 RLHF 中。多样性维护出现在进化计算、DGM、AlphaEvolve 和开放式科学中。奖励博弈出现在 RL、代码生成和基准驱动的智能体设计中。安全约束出现在任何代码或工具可被修改的地方。术语不同，但工程问题如出一辙。

统一还揭示了迁移机会。NAS 发展了权重共享和基准套件如 NAS-Bench；自我进化智能体需要类似的可复用评估景观。质量-多样性算法发展了存档和新颖性指标；代码进化智能体可以使用它们来避免局部最优。程序修复发展了测试生成和补丁验证方法；自我修改智能体需要它们来实现安全的代码进化。RL 发展了离策略评估和奖励模型诊断；自博弈和自验证方法需要它们来检测崩溃。软件工程发展了 CI/CD、沙箱和回滚；自我进化 AI 需要这些作为治理原语。

\subsubsection{贡献四：研究者网络映射}

第四项贡献是研究者网络图谱。该领域年轻，但文献揭示了三条汇合的谱系。进化计算谱系从 Risto Miikkulainen、Kenneth Stanley、Joel Lehman 和 Jeff Clune 延伸，经过 NEAT、新颖性搜索、质量多样性、POET、AI-GAs、ADAS、DGM 和 AlphaEvolve 相关工作 \cite{stanley2002evolving,lehman2011abandoning,clune2019ai,adas2025,darwingodelmachine2025}。语言智能体推理谱系经过 Karthik Narasimhan、Shunyu Yao、Noah Shinn，涵盖 ReAct、Tree of Thoughts、Reflexion、SWE-bench 和 SWE-agent \cite{yao2023react,yao2023tree,reflexion2023}。LLM 自我改进谱系包括 Aman Madaan、Peter Clark、Sean Welleck、Eric Zelikman，以及 Self-Refine、STaR、自奖励模型 \cite{selfrewarding2024} 和宪法式自我批评方面的工作 \cite{selfrefine2023,star2022,constitutionalai2022}。

这些谱系如今在机构和方法论上已相互重叠。Jeff Clune 的团队出现在 ADAS、DGM、AI Scientist 和开放式进化 AI 中。Shengran Hu 和 Cong Lu 连接了 ADAS 和 DGM；Cong Lu 也出现在 AI Scientist 相关工作中。Shunyu Yao 和 Karthik Narasimhan 连接了 ReAct、Reflexion、Tree of Thoughts、SWE-agent 和更广泛的语言智能体评估。AIWaves 和浙江大学贡献了 Agent Symbolic Learning，明确桥接了智能体框架和文本优化。Google DeepMind 贡献了 FunSearch 和 AlphaEvolve，将 LLM 引导的进化扩展到数学和基础设施。Microsoft Research 贡献了 ReVeal 的以验证为中心的代码智能体 RL。清华大学和 LeapLabTHU 贡献了 Absolute Zero 的零数据自博弈推理。

绘制网络图谱不仅仅是社会学分析。它解释了为什么某些想法会一起出现。DGM 的存档设计在 Clune 的开放式进化背景下是合理的。Reflexion 的基于记忆的语言 RL 契合普林斯顿/东北大学的语言智能体谱系。Agent Symbolic Learning 反映了产学研界对可部署智能体框架的兴趣。AlphaEvolve 反映了 DeepMind 的 AI for Science 和系统优化议程。AI Scientist 反映了 Sakana AI 的受自然启发的、开放式自动化关注点。理解这些谱系有助于研究者定位缺失的桥梁，例如形式化程序修复与智能体自我修改之间更强的联系，或持续学习理论与提示词/工具进化之间的联系。


\subsubsection{贯穿调查的运行示例与证据}

为使调查具体化，我们使用跨越五个循环的若干运行示例。这些示例的选择并非仅仅因为它们广为人知，而是因为每一个都揭示了自我进化 AI 中不同的设计决策。Self-Refine 展示了最小循环：无外部训练，无跨任务的显式记忆，无代码修改，仅在任务情节内进行迭代批评和修订 \cite{selfrefine2023}。Reflexion 添加了显式记忆缓冲区和评估器，将一次尝试的反馈转化为下一次的上下文 \cite{reflexion2023}。Agent Symbolic Learning 从情景改进走向组件优化，通过将智能体表示为具有提示词、工具和流水线节点的符号网络 \cite{agentsymboliclearning2024}。ADAS 随后将架构本身作为搜索对象，要求元智能体编写新的智能体程序 \cite{adas2025}。DGM 将该搜索扩展为自我改进智能体实现的开放式存档 \cite{darwingodelmachine2025}。AlphaEvolve 将同一广义思想扩展到工业算法发现和系统优化 \cite{alphaevolve2025}。Absolute Zero 和 ReVeal 展示了自我进化如何通过可验证奖励改变训练分布和策略权重 \cite{absolutezero2025,reveal2025}。

这些示例揭示，自我进化不是能力阶梯上的单一节点。在反馈主观且安全约束严格的场景中，提示词级循环可能比代码级循环更可靠。在编程任务中，代码级循环可能更强大，因为单元测试提供明确的反馈。权重级循环可能产生更深的泛化，但更慢、更昂贵、更容易受到奖励错误规约的影响。基于存档的循环可能是避免局部最优的最佳方式，但它使部署复杂化，因为系统不再具有单一规范状态。因此，本调查避免仅按表面复杂度排列方法。相反，它追问哪个循环适合哪个任务、验证器、成本机制和治理要求。

以 Self-Refine 作为第一个运行示例。其生成-批评-精炼结构足够简单，可以用单个模型和提示词模板实现。该方法报告了代码优化、数学推理、情感反转、对话、诗歌、可读性和约束生成方面的改进 \cite{selfrefine2023}。其重要性不在于任何一个基准数字，而在于证明基础模型可以将自身的自然语言反馈用作优化信号。但 Self-Refine 也暴露了一个核心弱点：产生错误的同一模型可能无法诊断该错误。没有外部评估器，自我批评可能变为自我合理化。在我们的框架中，Self-Refine 因此是循环 I 的一个高可访问性、低保证性的实例。

Reflexion 通过分离执行者、评估者和自我反思角色来强化循环 \cite{reflexion2023}。执行者尝试任务；评估者分配成功或失败；反思模型撰写语言化的经验教训；记忆将这些教训带入未来。在代码和交互式任务上，这种架构可以产生大幅增益，因为失败被转化为可复用的建议。Reflexion 报告的 HumanEval 结果以及在 AlfWorld 和 WebShop 上的改进，使语言强化学习成为核心范式。然而，Reflexion 也暴露了记忆管理问题。如果每次失败都成为记忆，上下文会膨胀；如果记忆被过度压缩，有用细节会丢失；如果反思是泛化的，它们增加词元数而不增加指导力。后续系统在存储文本反馈时都继承了这一权衡。

Agent Symbolic Learning 是我们循环 II 的主要贯穿示例，因为它将优化的对象从答案或记忆转变为智能体的符号化组件 \cite{agentsymboliclearning2024}。该方法的语言损失和语言梯度是对可微学习的类比，但这一类比在操作上是有用的。一条失败的轨迹可以被转换为文本诊断；诊断可以通过提示词、工具和流水线节点进行反向传播；然后每个节点可以由 LLM 优化器进行更新。在 HotPotQA、MATH、HumanEval、软件开发可执行性和创意写作上报告的增益表明，智能体组件可以以结构化的方式进行优化。与此同时，该方法也提出了新的问题：语言空间中的学习率是什么？语言梯度有多稳定？相互冲突的文本梯度应该如何聚合？如何防止提示词优化器在小规模开发集上过拟合？这些问题将符号化智能体优化与经典的优化理论联系起来，但并不能简化为后者。

ADAS 是我们自动智能体设计搜索的贯穿示例 \cite{adas2025}。它将实现智能体的 Python 程序视为搜索空间。这一点很重要，因为该空间不仅庞大，而且具有足够的表达能力来编码任意计算。元智能体可以发明新的分解过程、工具使用策略、记忆布局、投票方案、验证例程或递归调用。其前景是明确的：人类设计者不再需要枚举所有合理的智能体架构。其风险同样明确：生成的智能体可能是脆弱的、不安全的或难以解释的。ADAS 报告称，自动发现的智能体可以超越手工设计的基线，并在不同领域或模型骨干之间迁移。本综述使用 ADAS 来讨论何时对智能体代码的搜索优于手工工程，以及何时评估成本变得过高。

G{\"o}del Agent 是运行时自引用的贯穿示例 \cite{yin2024godel}。它不是维护一个独立后代群体，而是允许智能体通过猴子补丁（monkey patching）检查和修改自身的逻辑。这种设计接近于自我重写智能体的直觉概念。它也将软件工程的隐患推到了前台。运行时补丁可能引入不确定性、隐藏状态依赖、可移植性问题和回滚复杂性。该方法使用验证和回滚，但更广泛的教训是：随着自我进化从提示词走向可执行代码，系统必须采用来自安全软件部署的实践。补丁需要来源追溯、测试、沙箱、审计日志和明确的接受标准。

DGM 是我们开放式智能体进化的贯穿示例 \cite{darwingodelmachine2025}。其存档机制将自我改进从单条谱系转变为一个群体过程。这至关重要，因为一个智能体变体可能改进工具使用，另一个改进规划，第三个可能改进错误恢复。单一贪婪谱系可能丢弃那些不是即时最优但包含有用构建块的变体。存档保留了多样性并支持累积创新。DGM 在 SWE-bench 和 Polyglot 上报告的改进使其成为一个重要的实证里程碑。本综述使用 DGM 来分析存档采样、多样性奖励、跨领域迁移，以及改进智能体分数与改进智能体群体可进化性之间的区别。

AlphaEvolve 是我们具有强工业级评估的质量-多样性搜索的贯穿示例 \cite{alphaevolve2025}。它结合了多个 Gemini 模型、自动评估器和 MAP-Elites 风格的保留策略。其在矩阵乘法、数据中心调度、芯片设计和注意力核函数上报告的结果表明了为什么代码和算法是自我进化的吸引人领域：候选方案可以被执行、测量、比较和保留。然而 AlphaEvolve 也说明了基础设施方面的障碍。许多领域缺乏可靠、廉价的自动评估器。一个能在精确性能基准下将核函数改进 23\% 的系统，可能无法以同样的置信度改进政策论据、医疗决策或社交互动。因此，验证器的可用性是该领域的一个核心轴线。

Absolute Zero 是我们自生成课程学习的贯穿示例 \cite{absolutezero2025}。模型提出任务、求解任务，并从代码执行器获得验证。这种双重角色类似于博弈中的自我对弈，但以编程和推理任务取代了棋盘位置。这一范式之所以重要，是因为人类策划的数据可能成为进一步改进的瓶颈。如果模型能生成具有挑战性、可求解、多样化的任务，它就能创建自己的训练信号。但同样的设置也面临琐碎化、不可解任务、模式崩溃或课程短视的风险。除非奖励同时鼓励难度和可解性，否则提议者和求解者可能在任务空间的退化区域中共谋。因此，Absolute Zero 凸显了训练循环本身的治理问题。

ReVeal 是我们代码智能体中可靠自验证的贯穿示例 \cite{reveal2025}。它将验证不视为事后工具调用，而是视为应与生成一起优化的能力。这是一个重要的转变。许多智能体可以编写代码；但较少能设计有效的测试、解释失败并决定何时停止。ReVeal 的回合级信用分配和生成-验证协同进化指向了学会如何使用工具进行自我纠正的智能体。更广泛的教训是，自我进化不仅依赖于更好的生成器，还依赖于更好的评论者、验证器和信用分配机制。

AI Scientist 是生命周期级自我进化的贯穿示例 \cite{aiscientist2024}。它尝试自动化完整的研究周期：创意生成、新颖性检查、实验设计、代码执行、论文撰写和评审。它将综述的范围扩展到编码智能体和提示词优化器之外。在科学发现中，被改进的对象可能是假设、实验、论文或研究议程。反馈可能是实验结果、新颖性分数或同行评审标准。AI Scientist 也展示了当前系统的局限性：生成的论文可能缺乏理论深度，评审者可能不可靠，领域可能狭窄。这些局限性不是偶然的；它们表明长周期自我进化需要在复杂工作流的每个阶段都有可靠的中间评估。

这些贯穿示例共同构建了一架从局部修订到开放式系统重设计的阶梯。然而，这一阶梯不应被解读为一种编年史，即后来的方法简单地取代早期的方法。在实践中，强系统将组合多个循环。一个 DGM 风格的智能体可能在评估期间使用 Reflexion 式的记忆。一个 ADAS 生成的架构可能包含 Self-Refine 子程序。一个 ReVeal 训练的代码模型可能嵌入到 AlphaEvolve 式的群体搜索中。一个 AI Scientist 可能使用 Agent Symbolic Learning 来优化其评审提示词，并使用 DGM 式的存档来保留成功的实验策略。本综述的核心主张是：未来的领域将是组合式的——自我进化将是一个由交互循环组成的架构。

\subsubsection{为什么该领域现在需要一篇综述}

对综述的需求源于概念上的碎片化。同一种机制正在不同的标签下被重新发现：自我精炼、言语强化学习、递归内省、智能体符号学习、智能体系统自动设计、代码进化、自我对弈推理、可靠自验证、开放式发现和 AI 科学家自动化。没有共同的词汇，就很难比较各种声明。一篇论文可能报告一个智能体"自我改进"了，而实际上它只是在单个上下文窗口内修订了输出。另一篇可能声称"自我进化"，因为它在自生成的数据上训练，尽管部署的系统无法修改自身的组件。第三篇可能在群体中进化代码，但没有解决安全性或治理问题。本综述提供了一套词汇，可以精确表达：改变了什么、谁改变的、使用什么反馈、通过什么规则保留、在什么分布上评估。

该领域需要综述的另一个原因是实证结果难以比较。Self-Refine 报告了跨多种任务的平均改进；Reflexion 报告了代码和交互任务的增益；Agent Symbolic Learning 报告了在 QA、数学、代码、软件开发和创意写作上的改进；DGM 报告了 SWE-bench 和 Polyglot 的改进；AlphaEvolve 报告了科学和基础设施的发现；Absolute Zero 报告了零数据推理增益；ReVeal 报告了多轮代码智能体的改进。这些结果不共享基准套件、成本核算方法、安全协议或进化步骤的标准定义。静态模型排行榜无法捕捉一次性提示词精炼与 1000 次迭代存档搜索之间的差异。在社区能够构建公平的基准之前，综述级别的综合是必要的。

第三个原因是研究与开源采用之间的差距。流行的智能体框架通常宣传自主性、记忆或多智能体协作，但许多部署仍然是静态的。基于角色的 CrewAI 或 MetaGPT 工作流可以协调智能体而不修改自身。AutoGen 对话可能通过对话进行适应，但不一定保留结构性变更。LangGraph 可以表示动态工作流，但图的变更通常是因为开发者编辑了它。Letta/MemGPT 进化记忆，但不一定进化工具或代码。DSPy 编译提示词，但其自我进化取决于编译是否集成到反馈循环中。OpenEvolve 及相关项目将进化编码带入实践者，但验证器质量和基准选择仍然是决定性的。综述可以分离出使自我进化成为可能的基础设施与实际执行自我进化的系统。

第四个原因是安全性。自我进化之所以有吸引力，正是因为它可以发现人类未指定的变更。这也是它有风险的原因。一个自我修改的智能体可能移除约束、利用基准、学会满足代理指标而违反真正目标、毒害自身记忆，或生成带有隐藏漏洞的代码。一个自我对弈系统可能坍缩到琐碎任务中。一个反思系统可能存储错误的教训。一个科学自动化系统可能用看似合理但肤浅的论文淹没评审渠道。这些风险不能通过拒绝构建自我进化系统来解决——轻量级版本已经存在于提示词优化器、代码智能体和记忆系统中。社区需要一个共享框架来安全地设计它们。

第五个原因是自我进化改变了进步的度量单位。传统的单位是模型检查点。自我进化的单位是循环：一个不断进化的群体、一个提示词上的优化器、一个训练课程生成器，或一个修补自身的智能体。进步可能来自更好的验证器、更好的存档、更好的回滚、更好的记忆整合、更好的任务提议或更好的治理，而不仅仅是更大的基础模型。这种转变对研究策略很重要。如果未来的能力增益来自围绕基础模型自我改进的系统，那么基准报告必须包括进化协议、计算预算、评估器和保留规则，以及基础模型名称。

\subsubsection{开源证据和项目健康度的作用}

开源项目在本综述中扮演两个角色。首先，它们提供了实现，通过这些实现可以检查、复现和扩展自我进化机制。其次，它们塑造了更广泛社区理解自主性和智能体的语言。一个拥有数万星标的框架可以普及一种设计模式，即使其内部自我进化能力有限。相反，一个技术上重要的研究系统可能由于难以运行、成本高昂或与论文基准绑定而只有适度的采用。

本综述使用的星标分析证据警示不要将流行度等同于质量。炒作驱动的项目可以通过社交媒体、演示和自主性叙事快速积累星标。学术或工程密集型项目通常增长较慢，但吸引更高质量的分支、问题和贡献者。报告对 AutoGPT、Devika、CrewAI、MetaGPT、OpenEvolve、DSPy、SWE-Agent、OpenHands、LangGraph、AutoGen、FunSearch 及相关项目的比较促使了更审慎的评估。我们问的不是"哪个项目星标最多？"而是"哪个项目包含反馈驱动的修改循环，该循环有多可靠，有什么证据表明累积改进？"

这种区分对于将"Self Evolve"作为品牌和研究领域很重要。一个着陆页、索引或社区资源不应简单地聚合智能体项目。它应该识别每个项目在进化栈中的位置。AutoGPT 在普及自主任务循环方面具有历史重要性，但许多版本依赖于固定的提示词和工具使用。CrewAI 提供了可访问的基于角色的编排，但固定的团队不是自动自我进化的。MetaGPT 编码了软件开发角色，但其默认流水线很大程度上由人类设计。AutoGen 支持智能体对话和人在回路控制，这可以承载进化循环但不能保证它们。LangGraph 提供了灵活的状态图底层，可以表示自适应工作流。Letta/MemGPT 专注于记忆进化。DSPy 提供提示词/程序优化。SWE-Agent 和 OpenHands 将软件工程操作化，从代码仓库和测试中获取反馈。OpenEvolve 明确探索进化式代码改进。每一个都属于生态系统，但每一个都应按机制标注。

项目健康度也影响科学可信度。如果一个仓库已被废弃、缺乏文档或无法复现，那么其自我进化系统很难评估。分支质量、问题质量、贡献者多样性和拉取请求活动提供了间接证据，表明一个项目是否能支持累积研究。因此，综述将开源分析视为论文分析的补充。论文提供声明和受控实验；仓库揭示工程假设、工具依赖、评估脚本和社区采用。一个成熟的自我进化领域将需要两者兼备。

\subsubsection{综述中引入的评估原则}

由于自我进化系统随时间变化，其评估必须是时序性的。我们引入了几个在后续章节中反复出现的原则。第一个是\emph{轨迹报告}：论文不仅应报告最佳最终分数，还应报告分数序列、失败的修改、接受的修改、计算成本和评估方差。一个在数百个失败子代后找到一个强智能体的存档与一个逐步稳定改进的优化器是不同的。两者可能都有价值，但它们有不同的成本和风险特征。

第二个原则是\emph{留出进化}。一个系统不仅可能过拟合模型，还可能过拟合进化过程。如果修改在同一个基准上被反复选择，智能体可能学习到特定于基准的技巧。因此评估应区分用于进化的开发任务和用于最终评估的留出任务。DGM 的跨领域迁移声明很重要，因为它们解决了进化智能体是否能在搜索基准之外泛化的问题 \cite{darwingodelmachine2025}。类似的留出协议对提示词、记忆和课程进化也是需要的。

第三个原则是\emph{验证器审计}。在自我进化中，验证器是教师、裁判和选择环境。如果验证器是错误的，进化过程将优化向错误。代码执行器相对较强，但如果测试薄弱则仍不完整。LLM 评审者灵活但容易受到偏差和不一致性的影响。人类反馈有价值但昂贵且多变。奖励模型可以被博弈。验证器审计应包括对抗样本、校准检查、与人类的一致性，以及对提示词变化的敏感性。

第四个原则是\emph{成本归一化改进}。一种方法在数千次昂贵模型调用后改进一个百分点，可能不如一种以远低成本改进略少的方法有用。因此我们强调每次计算、每次令牌、每次环境步骤或每小时挂钟时间的改进。AlphaEvolve 规模的搜索对于数据中心调度或芯片设计可能是合理的，因为微小的百分比改进具有巨大价值 \cite{alphaevolve2025}。同样的成本对于低风险文本生成任务则不合理。

第五个原则是\emph{安全保留}。一个候选修改不应仅因为它改进了任务分数就被保留。它还应保留约束：安全性、隐私、工具权限、对齐指令、可复现性和回滚能力。这对代码级和架构级自我进化尤为重要。一个通过禁用测试来改进 SWE-bench 的补丁不是改进。一个通过忽略安全策略来增加用户满意度的提示词不是改进。接受函数必须同时编码能力和约束。

第六个原则是\emph{多样性和新颖性核算}。开放式系统可能发现不能立即改进主要指标的有用垫脚石。质量-多样性方法提供了一种保留此类变体的方式。然而，多样性不应成为保留不安全或不相关产物的借口。因此综述将多样性视为一个结构化目标，理想情况下由行为描述符或任务生态位定义，而非模糊的新颖性。

\subsubsection{安全与治理作为定义的一部分}

安全性在综述中通常被放在后面的章节，但对于 AI 自我进化而言，它属于定义部分。一个没有治理的自我进化系统不仅仅是一个不完整的产品；它是一个规格不良的优化器。接受函数必须决定什么算作改进。如果它只测量基准分数，系统可能牺牲可解释性、鲁棒性、隐私或合规性。如果它只测量用户认可，它可能变得谄媚。如果它只测量执行成功率，它可能利用测试。因此治理约束 $\mathcal{G}$ 是形式化系统元组的一部分。

治理始于范围。一个自我进化的智能体应该知道它被允许修改哪些组件。提示词编辑比工具权限变更风险更低；记忆更新比任意代码执行风险更低；本地沙箱补丁比生产部署风险更低。系统应暴露一个清晰的修改面。例如，一个智能体可能被允许编辑其任务特定提示词但不能编辑其安全前言。它可能被允许提议代码补丁，但在没有测试和人类批准的情况下不能应用。它可能被允许生成训练任务但不能更改奖励定义。范围限制将自我进化从不受控的能力转变为一种工程模式。

治理还需要来源追溯。每次修改都应可归因：什么反馈触发了它、什么模型提议了它、什么文件或组件改变了、运行了什么测试、什么分数改进或退步了、以及为什么它被接受。没有来源追溯，调试一个进化中的系统几乎是不可能的。Reflexion 记忆需要时间戳和任务链接；提示词优化器需要版本历史；代码进化智能体需要差异和测试日志；基于存档的系统需要父子谱系。这些记录不是官僚负担。它们是研究进化的实证基础。

回滚是另一个治理原语。一个自我进化系统应该能够撤销有害的修改。G{\"o}del Agent 在验证失败后明确包含回滚 \cite{yin2024godel}。软件工程实践提供了许多相关工具：版本控制、分阶段推出、金丝雀部署、持续集成、特性标志和审计轨迹。未来的自我进化智能体应默认集成这些实践。修改面越强大，回滚要求就越强。

人类监督仍然重要，但应放在有影响力的位置。要求人类批准每个小的反思就违背了低级自我改进的目的。对自我修改的生产代码不要求人类批准则不安全。分层方法更可取：低风险记忆条目的自动保留、提示词和工具编辑的验证门控保留、代码补丁的沙箱门控保留、以及影响权限、安全策略、外部副作用或部署的变更需要人类批准。这种分层治理反映了自我进化的层级。

最后，治理必须考虑多智能体设置。如果多个智能体可以修改共享的提示词、工具、记忆或代码，冲突和共谋就成为可能。一个智能体可能以损害另一个智能体的方式优化某个组件。一个评论者可能对生成器变得过于宽容。一个群体可能趋向于利用同一个验证器。因此多智能体自我进化需要所有权、锁定、审查和独立评估。这些关切在人类软件团队中已经很熟悉；当智能体参与自身改进时，它们变成了技术要求。

\subsubsection{后续章节的统一词汇}

在整个综述中，我们一致地使用几个术语。\emph{底层表示}（substrate）是被修改的表示：文本、记忆、提示词、工具模式、图、代码、任务分布、权重或群体。\emph{变异算子}（variation operator）提议对该底层表示的变更。\emph{评论者}（critic）解释性能并产生反馈。\emph{验证器}（verifier）提供更客观的信号，如测试或环境奖励。\emph{选择器}（selector）决定保留什么。\emph{存档}（archive）存储变体或历史。\emph{治理器}（governor）执行约束。\emph{进化轨迹}（evolution trace）是随时间变化的状态、反馈、修改和接受/拒绝决策的序列。

这些术语使我们能够比较那些表面看来不相关的系统。在 Self-Refine 中，底层表示是输出；变异算子是同一个 LLM；评论者是自反馈；选择器可能是一个停止规则。在 Reflexion 中，底层表示是记忆；评论者是反思模型；验证器是评估器；选择器将反馈追加到记忆中。在 Agent Symbolic Learning 中，底层表示是提示词/工具/流水线；变异算子是语言优化器；评论者是语言损失；选择器更新符号权重。在 ADAS 中，底层表示是智能体代码；变异算子是元智能体；验证器是基准；存档存储发现的智能体。在 DGM 中，存档是核心，父本选择成为进化算法的一部分。在 Absolute Zero 中，底层表示包括课程和权重；验证器是代码执行器；选择器是 RL 更新。

该词汇还有助于诊断缺失的组件。有变异但没有验证器的系统会漂移。有验证器但没有多样性机制的系统可能过拟合。有存档但没有治理的系统可能保留不安全的变体。有强评论者但弱选择器的系统可能积累嘈杂的记忆。有强选择器但弱变异的系统可能停滞。许多当前论文最好被理解为改进一个缺失组件：ReVeal 改进了验证；RAGEN 改进了轨迹级 RL 并识别了回声陷阱；DGM 改进了基于存档的代码进化；Agent Symbolic Learning 改进了符号信用分配；Absolute Zero 改进了无外部数据的课程生成。

\subsubsection{本综述范围的局限性}

尽管综述范围广泛，但它并不声称每个自适应 AI 系统都是自我进化的。我们聚焦于操作组件的修改足够明确以致可以分析的系统。这排除了大部分普通的监督学习、一次性提示词工程、静态智能体编排和仅人类参与的模型迭代。它也排除了缺乏已实现反馈循环的关于递归自我改进的推测性声明。我们的目标是构建已有机制的实证综述，而非科幻场景的目录。

我们也强调围绕 LLM 和智能体的系统，因为最近的领域集中在此。然而，自我进化本质上不局限于语言。机器人策略可以自我修改控制器，科学智能体可以进化实验协议，推荐系统可以进化目标和探索策略，多模态智能体可以适应感知-行动循环。形式化定义是表示无关的。LLM 的焦点反映了当前证据：语言模型作为组成 AI 系统的工件的变异算子异常强大。

另一个局限性是许多报告的结果是近期的、基于预印本的或难以复现的。一些系统不完全开源。一些基准数字可能随着评估改进而改变。一些开源项目统计数据具有时效性。因此我们将它们用作机制的证据，而非永久排名。持久的贡献是框架：可变底层表示、反馈、修改、验证、保留和治理。

最后，我们不假设自我进化总是可取的。对于许多应用，静态系统更安全、更便宜、更容易认证。一个医疗分诊模型、金融决策系统或法律助手可能需要严格的版本控制和人类批准。自我进化在任务多样、反馈可靠、改进可以在部署前被沙箱化的场景中最有吸引力。综述的目标不是推广不受控的自我修改，而是理解何时以及如何使反馈驱动的系统改进变得有用和安全。

\subsubsection{比较自我进化系统的设计维度}

本文综述的机制可以沿一组设计维度进行比较。第一个维度是\emph{自我访问}。一个系统可能无法访问其内部、只有只读访问、提议访问、沙箱写入访问或生产写入访问。提示词级别的方法通常只对本地上下文有写入访问。记忆智能体对记忆存储有写入访问。代码进化系统可以在沙箱中读写源文件。运行时自我修改的智能体可以补丁实时方法。权重级别的系统通过训练循环更新参数。自我访问的程度同时决定了能力和风险。一个只能追加反思的系统是受限的但相对安全；一个能重写自身工具权限的系统是强大的但危险的。

第二个维度是\emph{反馈客观性}。反馈从主观自评到外部验证的执行不等。Self-Refine 使用模型生成的批评，灵活但不总是可靠 \cite{selfrefine2023}。Reflexion 通常使用评估器或环境成功信号 \cite{reflexion2023}。SelfEvolve、ReVeal、Absolute Zero 和 AlphaEvolve 受益于可执行验证 \cite{selfevolve2023,reveal2025,absolutezero2025,alphaevolve2025}。AI Scientist 使用新颖性检查、实验指标和基于 LLM 的同行评审，混合了客观和主观信号 \cite{aiscientist2024}。反馈客观性不保证有用性——一个单元测试可能不完整——但它强烈影响改进是否能累积。

第三个维度是\emph{时间跨度}。一些循环在单次推理轮次内运行。其他在单个任务的重复尝试上运行、在基准套件上运行、在部署生命周期上运行或在智能体世代上运行。跨度越长，记忆整合、存档管理和漂移控制就越重要。单步批评可以通过最终答案质量来评估。一个运行一个月的进化智能体必须通过稳定性、迁移性、成本和安全事件来评估。长跨度还创造了延迟信用问题：一个今天看起来中性的修改可能在未来带来改进。

第四个维度是\emph{变更粒度}。输出编辑是细粒度且可逆的。提示词编辑可能影响许多未来任务。工具模式变更可以改变行动可供性。工作流图变更可以改变控制流。代码补丁可以引入任意新行为。权重更新可以以难以定位的方式改变潜在能力。群体更新改变可用智能体的分布。粒度决定了需要多少验证。细粒度变更可以通过局部检查来接受；粗粒度变更需要更广泛的回归测试套件。

第五个维度是\emph{保留拓扑}。一些系统覆盖单一状态。其他追加记忆。其他维护版本历史、分支或存档。贪婪覆盖简单但容易回归。只追加的记忆保留经验但可能膨胀或自相矛盾。版本化提示词允许回滚和消融实验。基于存档的进化，如 DGM 和 AlphaEvolve，保留了多条谱系 \cite{darwingodelmachine2025,alphaevolve2025}。群体结构对于开放式进化尤其重要，因为它们防止过早收敛，但它们也使部署选择更难：哪个变体应该用于哪个用户或任务？

第六个维度是\emph{信用分配}。当一个系统改进时，哪个组件导致了改进？在简单的 Self-Refine 循环中，答案可能是最新的批评。在多工具智能体中，改进可能来自提示词、工具调用、规划器、验证器或它们的交互。Agent Symbolic Learning 明确尝试将语言梯度分配到节点 \cite{agentsymboliclearning2024}。ReVeal 使用回合级自适应信用分配 \cite{reveal2025}。RAGEN 将轨迹分解为状态、思考、动作和奖励 \cite{wang2025ragen}。信用分配不仅是训练问题；它也是可解释性和治理问题。如果我们无法识别哪个变更有帮助，我们就无法自信地保留、迁移或审计它。

第七个维度是\emph{搜索压力}。一些系统使用弱压力，如单次自评。其他使用强压力，如在基准上的多代选择。强压力可以产生大幅增益，但也增加了对评估器的过拟合。这是进化计算和 AutoML 中的熟悉教训。当搜索过程强大时，基准成为训练环境的一部分。因此基准保密性、留出测试、对抗性评估和迁移检查变得更加重要。

第八个维度是\emph{人类角色}。人类可以提供目标、批准变更、标注反馈、设计验证器、检查日志或在失败后干预。一个有人类批准的系统如果提议并验证自身的修改，仍然可以是自我进化的。反过来，没有人类的系统不一定更好；它可能只是治理更少。因此我们将人类参与视为一个设计参数而非取消资格的因素。适当的水平取决于风险和可逆性。

第九个维度是\emph{部署耦合}。一些进化在研究沙箱中离线进行；一些在用户交互期间在线进行。离线进化允许受控评估和回滚。在线进化允许个性化和快速适应，但有将用户暴露给不稳定变体的风险。生产自我进化可能需要两阶段结构：在线日志生成候选修改，但接受在验证后离线进行。这种结构映射了软件工程中的持续集成，可以防止不安全的即时自我重写。

第十个维度是\emph{模型依赖}。许多自我进化循环依赖前沿 LLM 来保证变异质量。Agent Symbolic Learning 指出需要强骨干来计算语言梯度 \cite{agentsymboliclearning2024}。ADAS 和 DGM 依赖基础模型来提议有意义的智能体代码 \cite{adas2025,darwingodelmachine2025}。AlphaEvolve 使用双模型策略来兼顾广度和深度 \cite{alphaevolve2025}。一个只与最强专有模型配合的循环与一个能与更小开放模型配合的循环，在可复现性和成本属性上是不同的。未来研究应报告对基础模型强度的敏感性。

\subsubsection{进化循环的示例性形式分解}

比较系统的一种简洁方式是将每个自我进化轮次分解为六个函数：
\begin{equation}
(S_t, x_t) \xrightarrow{\mathrm{act}} \tau_t \xrightarrow{\mathrm{judge}} F_t \xrightarrow{\mathrm{diagnose}} G_t \xrightarrow{\mathrm{propose}} \Delta_t \xrightarrow{\mathrm{validate}} V_t \xrightarrow{\mathrm{retain}} S_{t+1}.
\end{equation}
这里 $\tau_t$ 是行为轨迹，$F_t$ 是反馈，$G_t$ 是诊断或类似梯度的解释，$\Delta_t$ 是候选修改，$V_t$ 是验证证据。不同的论文实例化或省略不同的函数。Self-Refine 将评判和诊断合并为自反馈，可能不使用单独的验证。Reflexion 将评估器反馈与反思性诊断分开。Agent Symbolic Learning 将诊断明确为语言梯度。SelfEvolve 使用解释器错误作为反馈，修订后的代码作为候选修改。DGM 使用基准评估作为验证，存档插入作为保留。ReVeal 训练验证步骤本身。Absolute Zero 使用任务提议和求解来同时生成 $x_t$ 和 $\tau_t$。

这种分解有助于识别哪里需要进步。如果行动薄弱，系统需要更好的基础模型或工具。如果评判薄弱，它需要更好的评估器。如果诊断薄弱，反馈不会转化为有用的变更。如果提议薄弱，修改在语法或语义上会很差。如果验证薄弱，有害的变更将被保留。如果保留薄弱，改进将丢失或回归将累积。一个系统只有其链条中最薄弱的函数那么强。

该分解还阐明了为什么语言模型异常有用。它们可以参与每个函数。它们可以通过解决任务来行动，通过批评输出来评判，通过解释失败来诊断，通过编辑提示词或代码来提议，通过生成测试来验证，通过总结教训来保留。但对所有函数使用同一模型可能产生相关错误。模型可能未能注意到自身的错误、产生讨好的诊断、提议一个满足自身有缺陷判断的补丁，然后保留该补丁。因此强系统会多样化角色：外部执行器、独立的评估器模型、人类审查、留出测试和存档级选择都可以减少相关失败。

也可以将进化循环视为近似的贝叶斯或进化推断。系统维护对有用修改的信念；反馈更新该信念；提议从高潜力区域采样；验证估计适应度；保留更新群体或状态。这一视角暗示了不确定性感知的自我进化。系统不是贪婪地接受最高分的修改，而是可能维护置信区间、探索不确定的变体，或对高风险变更要求更强的证据。此类方法在当前的 LLM 智能体工作中尚不成熟，该领域通常依赖于小规模基准样本和确定性的接受/拒绝规则。

\subsubsection{从基准提升到能力积累}

一个核心问题是，自我进化产生的是持久的能力积累，还是仅仅针对特定基准的提升。针对一个数据集优化的提示词可能在其他场景下失效。一个代码补丁可能通过公开测试但无法通过隐藏用例。一个自生成的课程可能只教授狭隘的技巧。一个归档可能积累的变体各自表现强劲但难以组合。因此，本综述区分了\emph{分数提升}与\emph{能力提升}。分数提升是在进化过程中使用的评估指标上衡量的。能力提升则迁移到预留任务、新环境或未来的进化步骤中。

能力积累需要可复用的结构。Reflexion 记忆只有在编码通用经验而非任务特定提示时才能积累。Agent Symbolic Learning 只有在提示词和工具更新改进了底层流程时才能积累。ADAS 只有在发现的架构能跨域或跨骨干网络迁移时才能积累。DGM 只有在归档变体为未来智能体提供垫脚石时才能积累。AlphaEvolve 只有在进化算法揭示出可跨问题族复用的模式时才能积累。Absolute Zero 只有在自生成任务能诱导出超越代码领域可迁移的推理技能时才能积累。AI Scientist 只有在生成的研究思路和实验策略能改进未来的科学工作而非仅产出孤立论文时才能积累。

这一区分也影响我们如何解读阴性结果。一个自我进化系统可能未能提升最终分数，但仍然发现了有用的组件、诊断方法或基准。反之，它可能提升了分数却损害了通用性。因此，后续章节提倡更丰富的报告：保留修改的消融实验、迁移测试、已发现策略的定性分析以及谱系图。目标不仅是知道 $S_T$ 优于 $S_0$，还要知道为什么更优以及这个原因在其他场景中是否重要。

能力积累的视角将自我进化与开放性（open-endedness）联系起来。在开放系统中，目标不仅是解决已知的基准，还要持续生成新的可能性。质量-多样性归档、新颖性度量、环境生成和自博弈课程都是保持发现能力的机制。挑战在于保持开放性的生产性。无约束的新颖性可能漂移到无关的产物中；纯粹的目标压力可能过早收敛。新颖性与实用性之间的平衡是从进化计算继承的核心设计问题之一。

\subsubsection{对构建自我进化系统的启示}

对于实践者而言，该框架建议采用分阶段的方法。最安全的起点通常是围绕固定智能体的局部反思或验证循环。例如，代码助手可以生成测试、运行测试并在响应前修改代码。研究助手可以对照原始文档审视摘要。客户支持智能体可以在记忆中存储经批准的修正。这些循环在保持修改范围狭窄的同时提供即时价值。

下一阶段是符号组件优化。提示词、检索查询、工具描述、路由策略和工作流图可以进行版本控制并针对验证任务进行调优。此阶段应使用显式的开发集、回归测试和回滚机制。这通常比权重微调更实际，因为这些组件可解释且编辑成本低。Agent Symbolic Learning 和类似 DSPy 的系统指向这个方向 \cite{agentsymboliclearning2024}。

进一步的阶段是沙盒化代码进化。智能体可以提出对其自身工具、规划器或评估器的补丁，但补丁应被隔离、测试和审查。此阶段受益于软件工程基础设施：单元测试、静态分析、依赖扫描、权限边界和持续集成。DGM 和 AlphaEvolve 展示了代码进化的优势，而 G{\"o}del Agent 则表明了运行时保障的必要性 \cite{darwingodelmachine2025,alphaevolve2025,yin2024godel}。

最先进的阶段是种群或权重级别的进化。在此阶段，系统生成任务、更新策略、维护归档或进化多个智能体。此阶段需要大量计算资源和强大的评估能力。当领域具有可靠的奖励或验证器且长期改进值得投入成本时，此阶段是合适的。Absolute Zero、RISE、RAGEN、ReVeal、DGM 和 AlphaEvolve 处于此阶段的不同位置 \cite{absolutezero2025,qu2024rise,wang2025ragen,reveal2025,darwingodelmachine2025,alphaevolve2025}。

在所有阶段中，实际建议都是相同的：使循环显式化。定义可变的基板。定义反馈。定义提议机制。定义验证。定义保留。定义回滚。定义日志。定义谁可以批准高风险变更。一个无法回答这些问题的自我进化系统不是工程系统；它是一个不受控制的实验。

\subsubsection{定义所开启的研究问题}

本章引入的定义引出了若干研究问题。第一，自我进化在什么最小条件下能改善泛化能力而非导致过拟合？这个问题需要理论和基准来建模在有限反馈上的重复选择。第二，语言梯度方法应如何形式化？自然语言诊断表达力强但难以聚合、组合或验证。第三，如何为正确性不可执行的领域设计验证器？许多有价值的任务涉及判断、创造力、社会背景或长期后果。

第四，智能体进化的归档应如何构建？DGM 和 AlphaEvolve 展示了保留多样性变体的价值，但归档大小、采样、剪枝、生态位定义和安全过滤仍然是开放问题。第五，自我进化系统如何避免奖励博弈和评估器利用？这个问题在强化学习中很常见，但当系统可以修改代码、提示词或测试时变得更加尖锐。第六，自生成课程如何保持挑战性和对齐性？Absolute Zero 提出了一条路径，但任务提议本身就是一个优化问题。

第七，如何认证部署后会变更的系统？传统认证假设一个固定的制品。自我进化需要认证一个过程，包括其约束和监控。第八，人类监督应如何分配？过多监督阻止有用的适应；过少监督则招致不安全的漂移。第九，自我进化与可解释性之间是什么关系？进化轨迹可能提供系统如何变化的解释，但权重级别的更新可能仍然是不透明的。第十，当多个自我进化智能体共享工具、记忆或环境时，它们如何交互？

这些问题推动了本综述的其余部分。它们也表明了为什么 AI 自我进化不是提示工程或 AutoML 的一个狭窄子话题。它是一个涵盖表示、算法、软件基础设施、评估科学和治理的系统问题。


一个最终的实际启示是，自我进化应以足够的细节报告，使其可以作为一个过程被复现。一篇论文应说明基础模型、提示词、工具、可变文件或组件、反馈来源、进化预算、接受阈值、回滚策略、随机种子、基准划分以及所有保留的变体。没有这些细节，读者只能观察最终制品，而必须猜测产生它的进化路径。这类似于报告一个训练好的模型却不提供数据或优化器。因为循环是研究对象，所以循环必须被文档化。因此，我们在整个综述中将进化轨迹、审计日志和谱系记录视为一等科学制品。

这一过程视角也改变了该领域讨论失败的方式。失败的修改不仅仅是需要隐藏的噪声；它们揭示了搜索空间、评估器弱点以及塑造未来设计的约束。一个安全失败并记录失败原因的系统可能比一个只报告精选成功的系统更具科学价值。为了使自我进化 AI 走向成熟，阴性证据、回归和被拒绝的补丁应与成功的改进一起被分享。

\subsection{综述方法论}
\label{subsec:intro-methodology}

本综述遵循从 PRISMA 范围综述指南改编的系统文献收集和分类协议 \cite{moher2009preferred}。该方法论包括四个阶段：来源识别、筛选、资格评估和分类。

\paragraph{来源识别。} 我们检索了五个电子数据库（arXiv、Semantic Scholar、Google Scholar、ACL Anthology 和 NeurIPS/ICML/ICLR 会议论文集），使用的查询字符串为：\texttt{("self-evolving" OR "self-improving" OR "self-modifying" OR "self-refining" OR "recursive self-improvement") AND ("AI agent" OR "LLM" OR "language model" OR "agentic system")}。检索涵盖 2022 年 1 月至 2026 年 5 月的出版物，捕捉了指令微调 LLM 广泛部署后该领域的兴起。我们还从纳入论文中进行了反向雪球抽样，并通过 Semantic Scholar 进行了正向引文追踪。为补充学术来源，我们收集了标记有智能体、自我进化或自我改进关键词的 486 个 GitHub 仓库，产生了 204 个带有模型卡片文档的分析项目。

\paragraph{筛选。} 在去除重复项后，我们根据以下纳入标准筛选了 327 篇候选论文：（1）系统修改其自身的一个或多个操作组件（提示词、记忆、工具、代码、架构、数据课程或权重）；（2）修改由来自交互、评估或验证的反馈驱动；（3）系统在多个回合中保留或积累修改；（4）有足够的技术细节可用于分类进化机制。我们排除了仅描述静态提示、使用固定数据集的常规微调或没有反馈驱动自我修改的智能体编排的论文。

\paragraph{资格与分类。} 符合条件的论文沿五个维度进行分类：可变组件（什么在进化）、反馈来源（什么驱动变化）、验证机制（什么验证变化）、保留拓扑（变化如何持续）和时间范围（进化何时发生）。每篇论文至少由作者进行两次独立分类，分歧通过讨论解决。由此产生的分类法形成了第~\ref{subsec:intro-contributions}~节中描述的五进化循环框架。我们进一步按自我进化覆盖范围对开源项目进行了分类：实现反馈驱动修改循环的项目被标记为"核心进化"，托管或启用此类循环的项目被标记为"进化基础设施"，提供智能体能力但没有反馈驱动修改的项目被标记为"静态编排"。

\paragraph{证据综合。} 定量证据从原始论文报告的基准结果中提取，关注评估协议、样本量、基础模型和可复现性。定性证据来自架构分析、代码检查和社区反馈（Hacker News、Reddit、X/Twitter 讨论）。我们不将基准分数视为永久排名；相反，我们将其作为特定条件下机制有效性的证据。所有分类数据和项目元数据维护在一个版本化仓库中，以支持更新和修正。

\subsection{与现有综述的关系}
\label{subsec:intro-survey-comparison}

若干综述涉及 AI 自我进化的某些方面，但没有一篇提供了本文所提供的跨领域统一和机制级别的分类法。我们沿四个维度区分我们的工作。

\paragraph{自我进化 LLM 智能体。} Wang 等人 \cite{wang2025selfevolving} 综述了自我进化的基于 LLM 的智能体，侧重于反馈驱动的改进循环。他们的分类法按进化阶段（经验获取、细化、更新）组织方法，而我们的五进化循环框架按反馈变为保留变化的机制进行分类。我们的综述还涵盖了他们以 LLM 为中心的范围所排除的进化计算、AutoML/NAS 和生产框架。CharlesQ9/Self-Evolving-Agents \cite{charlesq9selfevolving} 提供了 GitHub 索引的资源列表，但没有正式分析或跨领域比较。YoungDubbyDu/LLM-Agent-Optimization \cite{youngdubbydu2026optimization} 综述了 LLM 智能体的优化方法，侧重于提示词和流水线调优而非完整范围的可变组件。

\paragraph{AutoML 与神经架构搜索。} 经典 AutoML 综述 \cite{he2021automl,yao2018taking} 涵盖自动化流水线构建和超参数优化，但不涉及部署后的自我修改、智能体架构或代码级别的进化。NAS 综述 \cite{elsken2019neural,ren2021comprehensive} 侧重于神经网络拓扑搜索而非智能体软件、提示词或多模态架构。我们的综述通过将 ADAS \cite{adas2025} 和 DGM \cite{darwingodelmachine2025} 视为 NAS 应用于智能体程序的后继来桥接 AutoML/NAS 与智能体研究。

\paragraph{进化计算与 LLM。} 关于 LLM 增强进化算法的综述 \cite{yang2024large,liu2024large} 涵盖了 LLM 作为进化框架中的变异算子，但不涉及自指智能体系统、反思记忆或生产部署关注点。LLM4EC \cite{wuxingyu2025llm4ec} 侧重于 LLM 与进化计算在组合优化中的交叉。我们的综述将进化计算视为若干贡献传统之一，与强化学习、程序合成和智能体推理一起，统一在自我进化抽象之下。

\paragraph{LLM 智能体与推理。} 最近关于 LLM 智能体的综述 \cite{xi2025riseagent,wang2024surveyagent} 涵盖了工具使用、规划和多智能体协调，但按能力而非进化机制对方法进行分类。关于自我改进推理的综述 \cite{huang2025selfimprove} 涉及思维链细化但不涵盖代码级别或架构级别的自我修改。我们的综述的标志性贡献是跨所有七种可变组件类型（提示词、记忆、工具、代码、架构、数据、权重）比较方法的机制级别分类法，以及将安全、治理和生产部署作为自我进化定义的组成部分的显式处理。

\begin{table}[t]
\centering
{\TableTight
\begin{tabularx}{\textwidth}{@{}L{0.22\textwidth}cccc@{}}
\toprule
\textbf{维度} & \textbf{本综述} & \textbf{Wang 等人 2025} & \textbf{AutoML 综述} & \textbf{EC+LLM 综述} \\
\midrule
可变组件 & 7 种类型 & 3--4 种类型 & 1--2 种类型 & 1 种类型 \\
进化机制 & 5 个循环 & 基于阶段 & 基于流水线 & 基于算子 \\
部署后自我修改 & 是 & 部分 & 否 & 否 \\
智能体架构 & 是 & 是 & 否 & 否 \\
代码级进化 & 是 & 有限 & 否 & 是 \\
生产框架 & 是 & 否 & 否 & 否 \\
安全与治理 & 核心定义 & 单独章节 & 极少 & 极少 \\
跨领域统一 & AutoML+NAS+EC+RL+智能体 & 仅 LLM 智能体 & 仅 AutoML & 仅 EC \\
\bottomrule
\end{tabularx}}
\caption{本综述与相关综述在关键维度上的比较。本综述独特地涵盖了所有七种可变组件类型、五种进化循环机制、生产框架，并将安全视为自我进化定义的一部分。}
\label{tab:survey-comparison}
\end{table}

\subsection{论文结构}
\label{subsec:intro-structure}

综述的其余部分组织如下。第~2~章发展了 AI 自我进化的分类法。它扩展了这里引入的层次——输出、记忆、提示词/工具、架构、代码、课程、权重和种群——并按自主性、反馈来源、验证强度、保留机制和安全边界进行多维分类。该章还区分了闭环适应与开放进化，并引入了用于比较在不同层次上运行的系统的术语。

第~3~章研究核心方法。它涵盖反思性提示、言语强化学习、文本反向传播、提示词和工具优化、多智能体反馈、自调试以及验证引导的程序合成。它将 Self-Refine、Reflexion、Agent Symbolic Learning、SelfEvolve、ReVeal、RISE 和 RAGEN 视为方法族而非孤立的论文。该章强调反馈如何被表示并转化为修改。

第~4~章聚焦于进化和开放系统。它连接了 NEAT、新颖性搜索、质量多样性、POET、MAP-Elites、FunSearch、ADAS、DGM、AlphaEvolve、OpenEvolve 和 AI Scientist。该章解释了 LLM 如何改变变异算子的角色、归档如何支持累积发现，以及为什么开放性需要不同于单任务优化的评估原则。

第~5~章讨论评估。它论证了自我进化系统必须作为轨迹来评估。标准的 pass@k、准确率、胜率或奖励度量捕捉端点性能，但不捕捉可进化性、保留性、成本、安全性、迁移性或鲁棒性。该章综述了 HumanEval、MBPP、DS-1000、SWE-bench、LiveCodeBench、HotPotQA、MATH、ALFWorld、WebShop、ARC、GFootball 和 AI-scientist 同行评审设置等基准。它还讨论了评估器过拟合、基准泄漏、奖励博弈、回声陷阱和可复现性。

第~6~章综述框架和实现生态系统。它比较了智能体框架和开源项目，如 AutoGen、LangGraph、Letta/MemGPT、DSPy、SWE-Agent、OpenHands、CrewAI、MetaGPT、AutoGPT、smolagents、OpenEvolve、ADAS、DGM 实现、Self-Refine 仓库、Reflexion 实现和 AI Scientist。该章将静态编排与真正的自我进化区分开来，并使用星级数之外的信号分析项目健康度。

第~7~章研究痛点和风险。它涵盖不安全的自我修改、工具误用、沙盒逃逸、评估器博弈、奖励博弈、记忆投毒、提示词漂移、灾难性遗忘、成本爆炸、上下文窗口限制、来源追溯、问责制和治理。它还讨论了为什么形式化保证仍然困难，以及实际系统如何使用回滚、审计日志、分阶段部署、人工批准和受限动作空间。

第~8~章概述未来方向。它提出了自我进化轨迹的基准套件、进化历史的标准化日志格式、更安全的自修改代码协议、混合形式-经验验证、多智能体进化生态系统、自我进化的科学工作流以及可以改变自身行为的系统的治理标准。该章回归核心论题：未来 AI 进步不仅取决于更大的静态模型，还取决于 AI 系统改进自身的可靠机制。

总之，本引言将 AI 自我进化框架为一种实践和理论上的转变。该领域继承了图灵的行为评估、Schmidhuber 的自指雄心、AutoML 的自动化设计、进化计算的变异与选择、强化学习的反馈优化以及 LLM 革命操控语言和代码的能力。其定义性对象是修改系统本身的反馈驱动循环。因此，核心问题不再仅仅是"模型有多强？"或"智能体有多好？"，还包括"系统如何改变自身、谁验证变化、保留了什么、遗忘了什么、可能出现什么问题，以及我们如何衡量进化轨迹上的进步？"
